%! Composition of Functions — Unit 5, Lesson 5.6 — Warm-Up
\documentclass[10pt]{article}
\usepackage{algebra2-article}
\usepackage{algebra2-boxes}
\newcommand{\TallMath}[1]{$\displaystyle #1\rule[-1.4em]{0pt}{3.2em}$}

\begin{document}

\pageheader{Unit 5, Lesson 5.6}{Warm-Up}
\namedateperiod

\vspace{0.10in}

\begin{notesbox}{Getting Ready: Skills We'll Use Today}
\small Three short items. Item 1 part (d) is today's whole lesson --- you just have not been told
its name yet.

\vspace{4pt}
\textbf{1. Function notation (Algebra 1).} Let $f(x)=2x-3$ and $g(x)=x^{2}+1$.
\begin{center}
\renewcommand{\arraystretch}{1.7}
\begin{tabular}{@{}l l l@{}}
(a) \; $f(5)=$ \blank{1.6cm} &
(b) \; $g(5)=$ \blank{1.6cm} &
(c) \; $g(-1)=$ \blank{1.6cm} \\
\end{tabular}
\end{center}
\vspace{-6pt}
(d) Now take your answer to part (c) and feed it into $f$:
\[
  f\big(\text{your answer to (c)}\big)=f\big(\blank{1.2cm}\big)=\blank{1.4cm}
\]

\vspace{0.14in}

\textbf{2. Substituting a whole expression (Unit 3).} Let $h(x)=x^{2}-4x$. Replace \emph{every} $x$
with what is inside the parentheses, then simplify.
\begin{center}
\renewcommand{\arraystretch}{1.7}
\begin{tabular}{@{}l l l@{}}
(a) \; $h(3)=$ \blank{1.6cm} &
(b) \; $h(a)=$ \blank{2.6cm} &
(c) \; $h(x+1)=$ \blank{4.0cm} \\
\end{tabular}
\end{center}
\vspace{-4pt}
In part (c) the thing you substituted was not a number at all. What did you have to do to every
single $x$ in the rule? \writeline

\vspace{0.18in}

\textbf{3. Does order matter?} Start with the number $10$ and run two operations --- \emph{subtract
$4$} and \emph{double} --- in each of the two possible orders.
\begin{center}
\renewcommand{\arraystretch}{1.6}
\begin{tabularx}{0.96\linewidth}{@{}X c c@{}}
\hline
 & \textbf{After the first step} & \textbf{Final answer} \\
\hline
(a) Subtract $4$, \emph{then} double & \blank{1.8cm} & \blank{1.8cm} \\
(b) Double, \emph{then} subtract $4$ & \blank{1.8cm} & \blank{1.8cm} \\
\hline
\end{tabularx}
\end{center}
\vspace{-2pt}
Finish the sentence. \emph{Same two operations, same starting number --- but changing the}
\blank{3.0cm} \emph{changed the} \blank{3.0cm}.

\end{notesbox}

\end{document}
